\section{The Fruition, in Summary and Comparison}

The goal of the whole path is \textbf{buddhahood}, the full awakening of the mind to its own nature. It is not the making of something new. It is the recognition of what was always there.

\subsection{The three bodies}

A buddha is described by the \textbf{three bodies} (\textbf{three kayas}), present on three levels at once.

\begin{longtable}{L{3.2cm} L{8.0cm}}
\caption{The three bodies of a buddha.}\\
\toprule
\textbf{Kaya} & \textbf{What it is} \\
\midrule
dharmakaya & the truth body, formless unconditioned reality, the buddha-mind itself \\
sambhogakaya & the enjoyment body, subtle radiant form in pure realms, the peaceful and wrathful deities \\
nirmanakaya & the emanation body, physical form born in the world to teach \\
\bottomrule
\end{longtable}

The three correspond to mind, speech, and body, and to the three bardos at death. Some traditions add a fourth, the \textbf{unity of the three} (\textbf{svabhavikakaya}), and a fifth, the \textbf{unchanging continuity} running through them all (\textbf{vajrakaya}).

\subsection{The marks and qualities of a buddha}

A buddha is also described by classical sets of marks and powers. The body bears the \textbf{thirty-two major marks} (such as the urna, the curl between the brows, and the ushnisha, the crown protrusion) and the \textbf{eighty minor marks} that refine them. The awakened mind has the \textbf{ten powers} (unobstructed knowledge of karma, of the dispositions of beings, of past lives, of the destruction of the defilements, and more), the \textbf{four fearlessnesses}, and the \textbf{eighteen unique qualities} held by a buddha alone. The life of a supreme buddha unfolds in \textbf{twelve deeds}, from descent through enlightenment to passing into parinirvana.

\subsection{Samsara and nirvana are one reality}

The key point of the tradition is this. \textbf{Samsara and nirvana are not two separate places.} They are two ways of experiencing the same empty and luminous reality, one in ignorance and one in recognition.

The machinery of realms, hells, heavens, kalpas, and deities is the display of mind. To know mind is to complete the path. Nothing is escaped to. The same world, seen rightly, is the pure land. The same mind, recognizing its empty and luminous nature, is already a buddha.

\subsection{Comparison to the Kabbalah model}

The Kabbalah chapter and this Tibetan chapter are the two most fully worked-out systems in the book, and they are built on opposite foundations. Setting them side by side sharpens both.

\begin{longtable}{L{3.0cm} L{4.2cm} L{4.2cm}}
\caption{The Kabbalah and Tibetan models side by side.}\\
\toprule
\textbf{Question} & \textbf{Kabbalah} & \textbf{Tibetan Buddhism} \\
\midrule
the source & the Infinite, Ein Sof, a fullness from which all flows & no creator and no first cause, beginningless mind \\
how worlds arise & emanation, the Infinite contracts and pours out the worlds & dependent arising, mind spins out the worlds through ignorance and karma \\
the wound & the shattering of the vessels, the breaking that scatters the light & ignorance, mind failing to recognize its own nature \\
what is trapped & the divine spark caught in the shells & buddha-nature, the awakened mind obscured but never lost \\
the goal & tikkun and return, repairing creation and returning the light to the source & liberation, recognizing mind and seeing samsara as nirvana \\
\bottomrule
\end{longtable}

The contrast is clean. Kabbalah is a \textbf{theistic emanation} model. There is an Infinite, it overflows, the overflow breaks, light is trapped, and the work is to repair and return. Tibetan Buddhism is a \textbf{non-theistic dependent-arising} model. There is no Infinite and no overflow, mind simply fails to know itself, the worlds appear from that failure, and the work is to recognize and wake.

Yet the deep shape rhymes. Both place a \textbf{hidden perfection} at the center (the trapped spark, the obscured buddha-nature). Both diagnose a \textbf{primal fault} that veils it (the shattering, the ignorance). Both make the human task one of \textbf{recovery rather than acquisition}. Nothing is created in the cure. What was always there is uncovered.

\subsection{Tie to the universal theme}

This is the thread that runs through every model in the book. Under different names, in theistic and non-theistic dress, the same story keeps being told. \textbf{Something whole is veiled, and the work of a life is to wake up and return to the source.}

In Kabbalah the source is the Infinite, and one returns light to it by repair. In Tibetan Buddhism there is no source to return to, only a nature to recognize, yet the arc is the same. A fall into forgetting, a long path of remembering, and a homecoming that turns out to be a recognition of what was never actually left. The maps differ. The waking is the same.
